% !Mode:: "TeX:UTF-8"

\chapter{大语言模型生成过程的错误识别}

\section{引言}

大语言模型在复杂问答任务中已经表现出较强的推理和知识组织能力，但其生成过程仍然容易出现隐蔽错误。
这类错误不一定只体现在最终答案上，还可能出现在问题理解、检索式查询、中间信息生成或答案整合等环节。
在自检索问答场景中，模型会先通过显式思考判断是否需要补充知识，再生成类似检索查询的片段，
并基于自身参数化知识生成中间信息，最后给出答案。由于该过程不依赖真实外部检索结果，
中间信息本身可能包含事实偏差；同时，即使最终答案看似正确，推理链中也可能存在不一致或偶然命中的情况。
因此，仅用最终答案是否匹配标准答案来评价模型并不足以刻画问答系统的可信程度。

错误识别是提升模型可靠性的基础步骤。若系统能够判断某条生成轨迹是否存在过程错误，
并进一步指出错误发生的位置和原因，后续就可以利用该信息进行反思反馈、样本筛选或强化学习优化。
现有方法中，过程奖励模型（Process Reward Model, PRM）常被用于对推理步骤进行打分，
再通过低分步骤识别潜在错误。然而，PRM 分数通常是连续奖励信号，
如何将奖励分数稳定转化为错误标签仍然存在困难。另一类方法将大语言模型作为评判器
（LLM-as-a-judge），直接根据问题、模型完整输出和标准答案判断生成过程是否存在错误。
该方法能够生成自然语言解释，适合构造后续反思训练数据，但其效果依赖评判模型本身的判断能力。

本章围绕自检索问答中的错误识别问题，比较两类技术路线：
第一，使用 Qwen2.5-Math-PRM-7B 对长推理链切分后的片段进行过程奖励打分，
并基于 Z-score 异常检测识别错误片段；第二，使用基于 vLLM 部署的不同大语言模型作为评判器，
通过结构化提示词输出是否存在错误的二值判断和错误分析。
实验基于 SSRL 训练过程中收集的 rollout 数据构建评测集，
并使用 DeepSeek V4 Flash 进行自动标注。结果表明，PRM 异常检测具有较高精确率，
但召回率较低，容易漏检大量错误样本；相比之下，经过反思验证数据训练的 Qwen2.5-3B-Reflection
在准确率、召回率和 F1 值上取得更均衡的表现，更适合作为后续错误缓解和多模型协同验证的基础模块。

\section{背景知识和相关工作}

\subsection{自检索问答中的过程错误}

大语言模型可作为世界知识的模拟器。依据这一背景，研究使用自检索\cite{fan2025ssrlselfsearchreinforcementlearning}对模型参数化
知识的调用进行优化。自检索部分通过模型自身判断当前推理步骤是否需要对特定知识进
行检索，在确有必要时自动构建检索查询并针对查询依赖自身参数化知识进行知识调用，
无需外部查询的情况下用于搜索驱动型任务。相较于传统一次性使用检索的强化学习，自
检索策略能够降低检索带来的时间开销、避免冗余检索。相比普通思维链过程，自检索策
略能更好地对模型在事实部分掌握进行解析，也便于后续验证反思阶段对检索内容的合理性进行回
溯审查。

本文关注的模型输出并非只包含最终答案，而是包含完整的“思考-检索-回答”过程。
在该过程中，模型通常使用 \verb|<think>| 表示内部推理，
使用 \verb|<search>| 表示其认为需要查询的内容，
使用 \verb|<information>| 表示模型依据自身知识生成的中间信息，
最后使用 \verb|<answer>| 输出答案。与传统检索增强生成不同，
这里的 \verb|<information>| 并不来自外部知识库或搜索引擎，
而是由模型自行生成。因此，该过程可以展示模型如何调用和组织自身知识，
但也会将模型内部知识错误显式暴露出来。

自检索问答中的错误具有多样性。
模型可能在 \verb|<think>| 阶段错误理解问题，
也可能生成过宽、过窄或与问题无关的 \verb|<search>| 查询；
即使查询方向合理，模型在 \verb|<information>| 中仍可能生成错误事实；
最终答案还可能与前文信息不一致。上述错误可能单独出现，也可能在同一条轨迹中连续传播。
因此，错误识别任务需要从完整过程出发，而不能只检查最终答案字符串\cite{ref15_he2025large,ref16_you2025probabilistic}。

\subsection{基于过程奖励模型的错误检测}

过程奖励模型（Process Reward Model, PRM）是近年来大语言模型复杂推理增强与对齐领域的重要技术路线。
与针对最终输出结果进行全局打分的结果奖励模型（Outcome Reward Model, ORM）不同，
其原理是为推理链中每一步骤打分，以区分正确与错误步骤。
这一特性使其可以在训练中提供密集的监督信号，同时也可以对推理路径中每一步骤的逻辑连贯性与事实准确性进行量化评估，
从而区分正确与错误步骤。



对于一条较长的自检索轨迹，可以先将其切分为若干片段，
再由 PRM 对每个片段给出奖励分数\cite{ref15_he2025large}。若某个片段分数显著低于总体分布，
则可将其视为潜在错误片段。

本章采用的 PRM 基线沿用基于 Z-score 和全局阈值的判别方式。

\begin{align}
r_t &= R_{\mathrm{PRM}}(s_t \mid x, s_{<t}) 
\end{align}

式中 \(s_{<t} = (s_1, \ldots, s_{t-1})\) 表示在步骤 \(s_t\) 之前已经生成的推理部分，
\(r_t\) 表示过程奖励模型对步骤 \(s_t\) 的评分。


\begin{align}
t &= \mu - \sigma \\
r_t &< t \Rightarrow \mathrm{Error\ Section} 
\end{align}

式中 \(t\) 为判定推理步骤是否包含错误的临界阈值，
\(\mu\) 为过程奖励模型生成的数据集中奖励分布的均值，
\(\sigma\) 是数据集中奖励分布的标准差。
当 \(r_t < \mu - \sigma\)时，该片段被视为异常低分片段。
在样本级别，只要一条生成轨迹中存在至少一个异常低分片段，
系统就将该样本预测为存在过程错误。该方法实现简单，且不需要额外训练二分类器，能够识别出明显的错误片段。

\subsection{基于大语言模型评判器的错误识别}

LLM-as-a-judge 方法将错误识别转化为条件判断任务。
评判模型接收原始问题、待评估模型的完整输出和标准答案，
然后判断生成过程是否正确，并在存在错误时给出错误描述和反思反馈。
与 PRM 只输出片段分数不同，评判模型能够直接利用问题语义、标准答案和完整轨迹之间的关系，
从而更适合识别事实矛盾、查询偏移和答案整合错误\cite{ref15_he2025large}。

现有系统性研究\cite{zheng2023judging}表明，在提供明确评分细则和结构化提示词的前提下，先进的大语言
模型在复杂任务上的评判结果与人类专家的评判结果有着高一致性。

本章使用的评判提示词要求模型首先审查完整的“思考-检索-回答”过程，
随后说明错误原因或验证正确性，并以严格 JSON 格式输出，以被程序解析。
对于错误样本，模型应输出对于样本分析，是否错误的说明，对错误的描述，以及反思性质的反馈反馈可用于后续错误缓解章节中的反思训练提供文本监督。
对于正确样本，输出字段包括
样本分析，是否错误的说明，对错误的描述，以及验证正确性的分析。

\section{数据构建与方法设计}

\subsection{评测数据构建}

本章评测数据来自使用 Qwen2.5-3B-Instruct 进行自检索强化学习
（Self-Search Reinforcement Learning, SSRL）过程中生成的 rollout 样本。
每条样本包含用户问题、模型完整输出和标准答案。
模型完整输出保留了思考、检索式查询、中间信息和最终答案，
因而可以用于评估过程级错误，而不仅是最终答案是否正确。

本文使用 DeepSeek V4 Flash 对样本进行自动标注。
标注任务包括三部分：判断样本是否存在错误、分析错误原因以及对正确样本进行验证说明。
对于错误样本，标注结果还包含可用于修正模型输出的反思反馈。
共标注 10000 条数据，并按照 8:2 的比例划分训练集和测试集。
训练集用于训练或微调反思验证模型，测试集用于比较不同错误识别方法的泛化能力。
在本章实验中，DeepSeek V4 Flash 的标注结果被作为参考标签。

为检验自动标注结果的可靠性，本文从全部标注结果中随机抽取100条样本进行人工复核，检查反思反馈以及对于错误的分析正确。人工复核结果显示，自动标注与人工判断的一致率为 0.96，Kappa 系数为 0.8906，表明二者具有较高一致性。因此，本文将 DeepSeek V4 Flash 的自动标注结果作为后续错误识别实验的参考标签。

\subsection{PRM 异常检测方法}

PRM 方法使用 Qwen2.5-Math-PRM-7B 作为过程打分模型。
首先将长自检索轨迹切分为若干相对独立的片段，
再对每个片段计算过程奖励分数。随后，依据全局奖励分布计算 Z-score，
并使用 \(\mu-\sigma\) 对应的阈值识别异常低分片段。
如果一条样本中至少存在一个异常片段，则预测该样本存在错误；
否则预测该样本不存在错误。

该方法的优点是可解释性较强：被判定为错误的样本通常可以定位到低分片段。
但它也存在两个限制。
第一，固定阈值难以适配不同问题难度和不同轨迹长度，
同一分数在不同样本中的含义可能并不一致。
第二，PRM 的打分目标并非专门训练为错误二分类，
当错误需要结合标准答案或多个片段之间的一致性才能发现时，
单片段分数可能无法提供足够信号。

\subsection{评判模型方法}

评判模型方法使用大语言模型对错误进行识别检测。
大语言模型的输入包括问题、待检测的完整输出和标准答案，输出为可以被程序解析的 JSON 格式的错误分析和是否存在错误。
本章使用多个候选评判模型，包括 Qwen2.5-3B-Instruct、Qwen2.5-7B-Instruct、产生上述错误的 Qwen2.5-3B-SSRL
以及在反思验证数据上训练得到的 Qwen2.5-3B-Reflection。其中，Qwen2.5-3B-Instruct 和 Qwen2.5-7B-Instruct 用于观察通用指令模型的评判能力；
Qwen2.5-3B-SSRL 用于观察经过自检索强化学习后的生成模型是否自然具备自我判断能力；
Qwen2.5-3B-Reflection 则用于验证反思验证数据对错误识别能力的提升。

\section{实验}

\subsection{评价指标}

本章使用准确率、精确率、召回率和 F1 值评价错误识别结果。
本文将“推理过程中存在错误”定义为正类，
并定义 TP 表示真实存在错误且被预测为错误的样本数，
FP 表示真实无错误但被预测为错误的样本数，
TN 表示真实无错误且被预测为无错误的样本数，
FN 表示真实存在错误但被预测为无错误的样本数。
各指标定义如下：
\begin{align}
\mathrm{Accuracy} &= \frac{TP + TN}{TP + FP + TN + FN},\\
\mathrm{Precision} &= \frac{TP}{TP + FP},\\
\mathrm{Recall} &= \frac{TP}{TP + FN},\\
\mathrm{F1} &= \frac{2 \times \mathrm{Precision} \times \mathrm{Recall}}
{\mathrm{Precision} + \mathrm{Recall}}.
\end{align}


\subsection{PRM 错误检测结果}

使用 Qwen2.5-Math-PRM-7B 在本文所构建数据集进行异常检测后的主要指标如表~\ref{tab:error-detect-prm-main} 所示，
使用 Qwen2.5-Math-PRM-7B 的异常检测有明显的高精确率的同时存在低召回率。这表明大量真实错误样本
未被 PRM 给予明显低的分数。

\begin{table}[htbp]
\caption{PRM 异常检测方法的主要错误识别指标}
\label{tab:error-detect-prm-main}
\vspace{0.5em}\centering\wuhao
\begin{tabular}{lcccc}
\toprule
方法 & Accuracy & Precision & Recall & F1-Score \\
\midrule
Qwen2.5-Math-PRM-7B & 0.3799 & 0.9401 & 0.2055 & 0.3373 \\
\bottomrule
\end{tabular}
\end{table}

上述结果说明，单纯依赖异常低分片段进行错误识别较为保守。
PRM 能够捕捉到明显错误的推理片段，
但会相当一部分存在错误的片段不会被 PRM 给予低的分数。这意味着 PRM 并未有效识别出错误片段。
完整混淆矩阵和各项指标见附录表~\ref{tab:appendix-error-detect-prm}。

\subsection{评判模型错误识别结果}

不同评判模型在测试集上的三次运行平均结果如表~\ref{tab:error-detect-judge-main} 所示。
附录表~\ref{tab:appendix-error-detect-judge} 记录实验的具体数据。

\begin{table}[htbp]
\caption{不同评判模型的错误识别平均结果}
\label{tab:error-detect-judge-main}
\vspace{0.5em}\centering\wuhao
\begin{tabular}{lcccc}
\toprule
模型 & Accuracy & Precision & Recall & F1-Score \\
\midrule
Qwen2.5-3B-Reflection & 0.8318 & 0.8880 & 0.8937 & 0.8908 \\
Qwen2.5-3B-Instruct & 0.8203 & 0.8632 & 0.9101 & 0.8860 \\
Qwen2.5-7B-Instruct & 0.7824 & 0.8010 & 0.9535 & 0.8706 \\
Qwen2.5-3B-SSRL & 0.6625 & 0.9235 & 0.6111 & 0.7355 \\
\bottomrule
\end{tabular}
\end{table}

Qwen2.5-3B-Reflection 取得最高准确率和 F1 值，
其 Accuracy 为 0.8318，F1-Score 为 0.8908。
这说明经过反思验证数据训练后，小规模模型也能够获得较强的过程级错误识别能力。

Qwen2.5-3B-SSRL 总体表现较低。
这可能与模型对自己生成数据的评估过度自信有关。这一定程度上说明检测大语言模型推理过程中错误，不能仅靠生成错误样本模型更改提示词与上下文实现。

相比多个模型，Qwen2.5-3B-Reflection 在精确率和召回率之间取得更好平衡，
其结果表明，专门构造的反思验证数据能够有效提升模型的错误识别泛化能力。

\subsection{结果讨论}

PRM 与评判模型的对比揭示了两类方法的适用边界。
PRM 异常检测依赖各片段 PRM 模型的分数分布，可以作为错误识别中快速筛查方法的实现。
但由于召回率较低，难以作为唯一错误识别模块。

LLM-as-a-judge 方法能够较好进行错误识别。
从本章结果看，Qwen2.5-3B-Reflection 在保持较高召回率的同时也维持了较高精确率，
因此被选为后续错误缓解实验中的主要验证与反思反馈生成模型。

\section{本章小结}

本章研究了自检索问答场景中的大语言模型错误识别问题。
首先，本文构建了基于 SSRL rollout 的反思验证评测数据，
并使用 DeepSeek V4 Flash 对 10000 条样本进行错误标注、错误分析和验证说明。
随后，本文比较了 PRM 异常检测和 LLM-as-a-judge 两类方法。
实验结果表明，基于 Qwen2.5-Math-PRM-7B 的异常检测召回率低，不能识别出相当数量的真实错误。
基于评判模型的方法整体表现更好，
其中 Qwen2.5-3B-Reflection 取得 0.8318 的准确率和 0.8908 的 F1 值，在四个模型中有最好的表现。

这些结果说明，复杂问答中的过程错误往往需要结合完整上下文和标准答案进行整体判断。
仅依赖片段级奖励分数难以稳定识别事实错误、查询偏移和答案整合错误。
因此，后续章节将进一步利用反思验证模型生成的错误描述和修正反馈，
将错误识别结果用于提升模型在自检索问答任务中的可靠性。

% Local Variables:
% TeX-master: "../thesis"
% TeX-engine: xetex
% End:
